\chapter{Introduzione}
\label{cha:introduzione}

\section{Gli ammassi stellari}
\label{sec:ammassi-stellari}

Gli ammassi stellari sono sistemi gravitazionalmente legati,
costituiti da un insieme di stelle che interagiscono attraverso
la gravità. La loro evoluzione dipende sia dalla dinamica interna
sia dall'ambiente gravitazionale in cui si muovono
\cite{binneytremaine2008}.

Nel corso del tempo, le interazioni gravitazionali e l'influenza
del campo galattico possono modificare la distribuzione delle
stelle e provocare una perdita progressiva di massa. Comprendere
questi processi richiede di studiare come le condizioni iniziali
del sistema e le caratteristiche della sua orbita ne influenzino
l'evoluzione dinamica \cite{computationalmethods2023}.

Le simulazioni numeriche a molti corpi, o simulazioni
\textit{N-body}, permettono di seguire il moto delle singole
particelle e di ricostruire l'evoluzione del sistema nel tempo.
Esse costituiscono quindi uno strumento per analizzare gli
effetti delle interazioni gravitazionali e del campo esterno
sulla struttura di un ammasso stellare.

\section{Galactic Tidal Stripping}
\label{sec:galactic-tidal-stripping}

Un ammasso stellare che orbita all'interno di una galassia
è soggetto al campo gravitazionale galattico. La differenza
di accelerazione gravitazionale tra le diverse regioni
dell'ammasso produce effetti mareali, che possono favorire
la perdita di stelle e modificare la struttura del sistema.
Questo processo è noto come \textit{Galactic Tidal Stripping}.

L'intensità degli effetti mareali dipende, tra gli altri
fattori, dalla posizione dell'ammasso rispetto al centro
galattico e dalle caratteristiche della sua orbita.
In particolare, orbite differenti possono esporre l'ammasso
a condizioni gravitazionali differenti, influenzandone
l'evoluzione e la perdita di massa
\cite{baumgardt2003}.

Lo studio di questo fenomeno richiede pertanto di considerare
contemporaneamente la dinamica interna dell'ammasso e il suo
moto nel potenziale gravitazionale esterno. La complessità
del problema rende utile l'impiego di simulazioni numeriche,
che permettono di esaminare l'evoluzione del sistema per
diverse condizioni iniziali.

\section{Obiettivo}
\label{sec:obiettivo}

